% Results by research question (Framework section). Tables from scripts/make_rq_tables.py; paired counts from
% results/analysis_tuned/paired.csv (item bootstrap within each window, 90% intervals).
\newcommand{\rqAvgBestCoarse}{8}
\newcommand{\rqAvgBestItem}{9}
\newcommand{\rqAvgBestBlock}{5}
\newcommand{\rqAvgBestCoarseWord}{eight}
\newcommand{\rqAvgBestItemWord}{nine}
\newcommand{\rqAvgBestBlockWord}{five}
\newcommand{\rqCells}{27}
\newcommand{\rqAvgBestCells}{22}
\newcommand{\rqAvgWorstCells}{0}
\newcommand{\rqWmWorstCells}{16}
\newcommand{\rqLgbWorstCells}{11}
\newcommand{\rqAvgSecondBlock}{4}
\newcommand{\rqBlockMfiveWm}{0.598}
\newcommand{\rqBlockMfiveLgb}{0.953}
\newcommand{\rqBlockMfiveAvg}{0.589}
\newcommand{\rqBlockFavWm}{0.482}
\newcommand{\rqBlockFavLgb}{0.565}
\newcommand{\rqBlockFavAvg}{0.412}
\newcommand{\rqBlockSdrWm}{0.576}
\newcommand{\rqBlockSdrLgb}{0.632}
\newcommand{\rqBlockSdrAvg}{0.470}
\newcommand{\rqPAvgWmBetter}{128}
\newcommand{\rqPAvgWmWorse}{27}
\newcommand{\rqPAvgWmN}{312}
\newcommand{\rqPAvgWmCoarsebetter}{68}
\newcommand{\rqPAvgWmCoarseworse}{19}
\newcommand{\rqPAvgWmCoarsen}{212}
\newcommand{\rqPAvgWmItembetter}{60}
\newcommand{\rqPAvgWmItemworse}{8}
\newcommand{\rqPAvgWmItemn}{100}
\newcommand{\rqPAvgLgbBetter}{124}
\newcommand{\rqPAvgLgbWorse}{21}
\newcommand{\rqPAvgLgbN}{312}
\newcommand{\rqPAvgLgbCoarsebetter}{77}
\newcommand{\rqPAvgLgbCoarseworse}{17}
\newcommand{\rqPAvgLgbCoarsen}{212}
\newcommand{\rqPAvgLgbItembetter}{47}
\newcommand{\rqPAvgLgbItemworse}{4}
\newcommand{\rqPAvgLgbItemn}{100}
\newcommand{\rqPWmLgbBetter}{64}
\newcommand{\rqPWmLgbWorse}{75}
\newcommand{\rqPWmLgbN}{312}
\newcommand{\rqPWmLgbCoarsebetter}{40}
\newcommand{\rqPWmLgbCoarseworse}{45}
\newcommand{\rqPWmLgbCoarsen}{212}
\newcommand{\rqPWmLgbItembetter}{24}
\newcommand{\rqPWmLgbItemworse}{30}
\newcommand{\rqPWmLgbItemn}{100}
\newcommand{\rqAvgItemWorseMax}{8}
\newcommand{\rqMcsWorldModel}{35}
\newcommand{\rqMcsLightgbm}{37}
\newcommand{\rqMcsAverage}{51}
\newcommand{\rqMcsEnsemble}{41}
\newcommand{\rqMcsTimesfm}{12}
\newcommand{\rqMcsCells}{60}
\newcommand{\rqFastShareDomLo}{11}
\newcommand{\rqFastShareDomHi}{21}
\newcommand{\rqFastShareRestLo}{37}
\newcommand{\rqFastShareRestHi}{42}
\newcommand{\rqShortShareDomLo}{11}
\newcommand{\rqShortShareDomHi}{21}
\newcommand{\rqShortShareRestMax}{2}
\newcommand{\rqAvgBestFast}{eight}
\newcommand{\rqAvgBestSlow}{eight}
\newcommand{\rqShortWmBeatsLgb}{eight}
\newcommand{\rqShortFavWm}{0.539}
\newcommand{\rqShortFavLgb}{0.667}
\newcommand{\rqShortMfiveWm}{0.866}
\newcommand{\rqShortMfiveLgb}{1.024}
\newcommand{\rqNineWord}{nine}
\newcommand{\rqSnaWm}{1.009}
\newcommand{\rqSnaLgb}{0.865}
\newcommand{\rqSnaAvg}{0.884}
\newcommand{\rqRuleLevel}{eight}
\newcommand{\rqRuleSna}{0.970}
\newcommand{\rqCostWmItemLo}{0.008}
\newcommand{\rqCostWmItemHi}{0.040}
\newcommand{\rqCostLgbItemLo}{0.058}
\newcommand{\rqCostLgbItemHi}{0.311}
\newcommand{\rqFaWmItemLo}{89}
\newcommand{\rqFaWmItemHi}{104}
\newcommand{\rqFaLgbItemLo}{66}
\newcommand{\rqFaLgbItemHi}{93}
\newcommand{\rqFaAvgItemLo}{78}
\newcommand{\rqFaAvgItemHi}{96}
\newcommand{\rqWmFaItemOff}{11}
\newcommand{\rqKeptSets}{five}
\newcommand{\rqItemWmBeatsAvg}{five}
\newcommand{\rqStoreWmBeatsLgb}{four}
\newcommand{\rqStoreAvgBeatsWm}{one}
\newcommand{\rqStoreWmBeatsAvg}{one}
\newcommand{\rqStoreLevel}{three}
\newcommand{\rqLgbFaItemLoFrac}{two thirds}
\newcommand{\rqCarriedWm}{0.447}
\newcommand{\rqOwnWm}{0.428}
\newcommand{\rqCarriedLgb}{0.450}
\newcommand{\rqOwnLgb}{0.448}
\newcommand{\rqCarriedWmIxs}{0.002}
\newcommand{\rqCarriedLgbIxs}{0.001}
\newcommand{\rqPoolGainMax}{0.007}
\newcommand{\rqPoolGainWm}{0.005}
\newcommand{\rqPoolGainWmSet}{analgesics}
\newcommand{\rqPoolGainLgb}{0.007}
\newcommand{\rqPoolGainLgbSet}{frozen entrees}
\newcommand{\rqPoolLossWm}{0.141}
\newcommand{\rqPoolLossWmSet}{soft drinks}
\newcommand{\rqPoolLossLgb}{0.235}
\newcommand{\rqPoolLossLgbSet}{analgesics}
\newcommand{\rqPoolGainMaxTwo}{0.01}


The results follow the four research questions. Every number comes from the
tuned world model and the tuned LightGBM, each set by the same search on each
data set (Section~\ref{sec:loop}), scored on the paper's windows. The average is
the equal-weight mean of the two forecasts.

\subsection{RQ1: grain}
\label{sec:rq-grain}

\textbf{Use the average at every grain.} Across the nine data sets and the
three views of Table~\ref{tab:rq-grain} (store grain and above by period, item
grain by period, and the whole window as one block), the average has the lowest
error in 22 of 27 cases and is never the least accurate of the three. The world
model is the least accurate in 16 cases and LightGBM in 11.

The paired tests give the same answer. They compare two forecasts level by level
and window by window, and count a difference only when it holds after
resampling items. Against the world model, the average is clearly more accurate
in 128 of 312 comparisons and clearly less accurate in 27; against LightGBM, 124
and 21. At item grain it is clearly less accurate than either model in at most 8
of 100 comparisons. Across data sets and levels it is in the 90\% model
confidence set in \rqMcsAverage{} of \rqMcsCells{} cases, against
\rqMcsLightgbm{} for LightGBM and \rqMcsWorldModel{} for the world model.

\textbf{Grain does not decide between the two models on their own.} The world
model is clearly more accurate than LightGBM in 40 comparisons at store grain
and above and clearly less accurate in 45; at item grain the counts are 24 and
30. Neither is ahead at either grain across the data sets, so a retailer that
picks one model for a grain is picking by data set, not by grain. The average
removes that choice.

\textbf{On the whole promotion block, the average is again the safe choice.} It
is lowest on five of nine data sets and second on the other four. LightGBM is
lowest on analgesics, bottled juice and snacks, but drifts in total where the
world model does not: its block error on M5 is 0.953, against 0.598 for the world
model and 0.589 for the average, and the same pattern shows on Favorita and soft
drinks.

\begin{table*}[t]
\centering\footnotesize\setlength{\tabcolsep}{3pt}
\caption{RQ1, grain. WRMSSE of the tuned world model (WM), the tuned LightGBM (LGB) and their equal-weight average
(Avg), averaged over the levels at store grain and above, over the item levels, and over every level for the whole
forecast window summed as one block; mean over the paper's windows. Lowest of the three in bold.}
\label{tab:rq-grain}
\resizebox{\textwidth}{!}{\begin{tabular}{l ccc ccc ccc}
\toprule
 & \multicolumn{3}{c}{Store and above, by period} & \multicolumn{3}{c}{Item levels, by period} & \multicolumn{3}{c}{Whole block, all levels}\\
\cmidrule(lr){2-4}\cmidrule(lr){5-7}\cmidrule(lr){8-10}
Data set & WM & LGB & Avg & WM & LGB & Avg & WM & LGB & Avg\\
\midrule
Analgesics & 0.688 & 0.698 & \textbf{0.678} & 0.710 & 0.705 & \textbf{0.695} & 0.523 & \textbf{0.458} & 0.472 \\
Bottled juice & 0.569 & 0.566 & \textbf{0.539} & 0.581 & 0.596 & \textbf{0.563} & 0.545 & \textbf{0.404} & 0.445 \\
Cheese & 0.381 & 0.471 & \textbf{0.377} & 0.623 & 0.594 & \textbf{0.578} & 0.445 & 0.425 & \textbf{0.387} \\
Frozen entrees & 0.339 & 0.327 & \textbf{0.320} & 0.375 & 0.350 & \textbf{0.344} & 0.353 & 0.330 & \textbf{0.318} \\
Soft drinks & 0.822 & 0.805 & \textbf{0.754} & 0.589 & 0.555 & \textbf{0.547} & 0.576 & 0.632 & \textbf{0.470} \\
Snacks & 1.241 & \textbf{1.040} & 1.071 & 0.545 & 0.516 & \textbf{0.508} & 0.893 & \textbf{0.633} & 0.688 \\
Frat & 0.471 & 0.479 & \textbf{0.449} & 0.528 & 0.524 & \textbf{0.508} & \textbf{0.345} & 0.406 & 0.364 \\
M5 & 0.660 & 0.634 & \textbf{0.591} & 0.859 & 0.886 & \textbf{0.856} & 0.598 & 0.953 & \textbf{0.589} \\
Favorita & 0.314 & 0.325 & \textbf{0.278} & 0.655 & 0.696 & \textbf{0.653} & 0.482 & 0.565 & \textbf{0.412} \\
\bottomrule
\end{tabular}}
\end{table*}


\subsection{RQ2: nature of the line}
\label{sec:rq-lines}

\textbf{Sales are concentrated, and the average holds up on both the fast lines
and the slow tail; where history is short, the world model is ahead of
LightGBM.} On the six Dominick's categories, 11 to 21\% of the items make 80\% of
unit sales over the year before the forecast; on M5, Frat and Favorita the head
is flatter, at 37 to 42\% (Table~\ref{tab:rq-lines}). On the fast lines the
average has the lowest error on eight of nine data sets, and on the slow lines
also on eight; on snacks LightGBM is lowest on the fast lines, and on cheese the
world model is lowest on the slow lines.

On series whose first sale falls less than 26 weeks before the forecast, the
world model's error is lower than LightGBM's on eight of nine data sets, Frat
being the exception: for example 0.539 against 0.667 on Favorita and 0.866
against 1.024 on M5. These series are 11 to 21\% of the series on the Dominick's
categories, where items enter and leave the range over the years, and under 2\%
on M5, Frat and Favorita.

\begin{table*}[t]
\centering\footnotesize\setlength{\tabcolsep}{3pt}
\caption{RQ2, nature of the line. Share of items that make 80\% of unit sales over the year before the forecast, and
WRMSSE averaged over every level, by period, on the fast lines (those items), the slow lines (the other items) and the
series with short history (first sale less than 26 weeks or 182 days before the forecast). WM, LGB and Avg as in
Table~\ref{tab:rq-grain}. Shading marks the lowest of the three: \colorbox{wmcell}{world model}, \colorbox{lgbcell}{LightGBM}, \colorbox{avgcell}{average}.}
\label{tab:rq-lines}
\resizebox{\textwidth}{!}{\begin{tabular}{l c ccc ccc ccc}
\toprule
 & Items making & \multicolumn{3}{c}{Fast lines} & \multicolumn{3}{c}{Slow lines} & \multicolumn{3}{c}{Short history}\\
\cmidrule(lr){3-5}\cmidrule(lr){6-8}\cmidrule(lr){9-11}
Data set & 80\% of units & WM & LGB & Avg & WM & LGB & Avg & WM & LGB & Avg\\
\midrule
Analgesics & 17\% & 0.754 & 0.748 & \cellcolor{avgcell}\textbf{0.738} & 0.756 & 0.769 & \cellcolor{avgcell}\textbf{0.748} & \cellcolor{wmcell}\textbf{1.291} & 1.459 & 1.350 \\
Bottled juice & 21\% & 0.594 & 0.606 & \cellcolor{avgcell}\textbf{0.571} & 0.513 & 0.516 & \cellcolor{avgcell}\textbf{0.494} & \cellcolor{wmcell}\textbf{1.303} & 1.393 & 1.312 \\
Cheese & 17\% & 0.505 & 0.507 & \cellcolor{avgcell}\textbf{0.461} & \cellcolor{wmcell}\textbf{0.515} & 0.584 & 0.526 & \cellcolor{wmcell}\textbf{1.222} & 1.411 & 1.284 \\
Frozen entrees & 19\% & 0.327 & 0.310 & \cellcolor{avgcell}\textbf{0.306} & 0.461 & 0.412 & \cellcolor{avgcell}\textbf{0.402} & 1.473 & 1.471 & \cellcolor{avgcell}\textbf{1.336} \\
Soft drinks & 11\% & 0.640 & 0.602 & \cellcolor{avgcell}\textbf{0.583} & 0.604 & 0.651 & \cellcolor{avgcell}\textbf{0.595} & 1.379 & 1.397 & \cellcolor{avgcell}\textbf{1.334} \\
Snacks & 17\% & 0.650 & \cellcolor{lgbcell}\textbf{0.496} & 0.524 & 1.189 & 1.154 & \cellcolor{avgcell}\textbf{1.144} & \cellcolor{wmcell}\textbf{4.204} & 4.565 & 4.328 \\
Frat & 42\% & 0.432 & 0.442 & \cellcolor{avgcell}\textbf{0.414} & 0.659 & 0.626 & \cellcolor{avgcell}\textbf{0.622} & 2.703 & 2.642 & \cellcolor{avgcell}\textbf{2.608} \\
M5 & 37\% & 0.717 & 0.700 & \cellcolor{avgcell}\textbf{0.662} & 0.735 & 0.748 & \cellcolor{avgcell}\textbf{0.702} & \cellcolor{wmcell}\textbf{0.866} & 1.024 & 0.907 \\
Favorita & 42\% & 0.417 & 0.442 & \cellcolor{avgcell}\textbf{0.395} & 0.496 & 0.522 & \cellcolor{avgcell}\textbf{0.472} & 0.539 & 0.667 & \cellcolor{avgcell}\textbf{0.525} \\
\bottomrule
\end{tabular}}
\end{table*}


\subsection{RQ3: seen and unseen data}
\label{sec:rq-unseen}

\textbf{On items it was not trained on, the world model loses almost no accuracy
and LightGBM loses a great deal; on stores it was not trained on, the world model
loses less on three of five data sets and the same on two.} With a fifth of the
items kept out of training, the world model's error on them rises by 0.008 to
0.040, and LightGBM's by 0.058 to 0.313; the difference is clear on all five data
sets tested (Table~\ref{tab:rq-transfer}). LightGBM under-forecasts the kept-out
items, at 66 to 93\% of actual sales; the world model forecasts 89 to 104\%. With a
fifth of the stores kept out, the world model loses less on cheese, snacks and
Frat, and the two lose the same on M5 and Favorita.

\begin{table*}[t]
\centering\footnotesize\setlength{\tabcolsep}{3pt}
\caption{RQ3, seen and unseen data. A fifth of the items, or of the stores, is kept out of training; the cost is the
WRMSSE on that fifth minus the WRMSSE on the same fifth of the same model trained on everything, mean over windows.
The interval is a 90\% bootstrap interval over windows for the world model's cost minus LightGBM's; ``clearly lower''
names the model whose cost is lower when the interval excludes zero. Forecast over actual is on the kept-out series.}
\label{tab:rq-transfer}
\resizebox{\textwidth}{!}{\begin{tabular}{ll cc c l cc}
\toprule
 & & \multicolumn{2}{c}{Cost of not training on them} & & & \multicolumn{2}{c}{Forecast over actual}\\
\cmidrule(lr){3-4}\cmidrule(lr){7-8}
Data set & Kept out & WM & LGB & WM $-$ LGB, 90\% & Clearly lower & WM & LGB\\
\midrule
Cheese & Items & 0.011 & 0.058 & [-0.064, -0.030] & World model & 0.99 & 0.93 \\
Cheese & Stores & 0.001 & 0.047 & [-0.060, -0.032] & World model & 1.05 & 0.94 \\
Snacks & Items & 0.012 & 0.162 & [-0.195, -0.109] & World model & 0.89 & 0.68 \\
Snacks & Stores & 0.001 & 0.103 & [-0.136, -0.068] & World model & 1.00 & 0.72 \\
Frat & Items & 0.008 & 0.313 & [-0.398, -0.248] & World model & 0.95 & 0.66 \\
Frat & Stores & 0.007 & 0.206 & [-0.210, -0.187] & World model & 0.99 & 0.77 \\
M5 & Items & 0.040 & 0.083 & [-0.057, -0.029] & World model & 1.04 & 0.84 \\
M5 & Stores & 0.022 & 0.020 & [-0.001, +0.006] & Neither & 0.99 & 0.93 \\
Favorita & Items & 0.010 & 0.171 & [-0.172, -0.145] & World model & 1.01 & 0.77 \\
Favorita & Stores & 0.020 & 0.017 & [-0.005, +0.015] & Neither & 0.97 & 0.98 \\
\bottomrule
\end{tabular}}
\end{table*}


\subsection{RQ4: transferability}
\label{sec:rq-transfer}

\textbf{The settings carry to another retailer, and training on other retailers'
data does not help either model.} M5's tuned settings carried unchanged to
Favorita give 0.447 overall for the world model, against 0.428 with Favorita's own
settings, and 0.450 for LightGBM, against 0.448 (Table~\ref{tab:rq-retailer});
at item by store the difference is 0.002 for the world model and 0.001 for
LightGBM. Training each model on its data set together with other data sets
lowers the error by no more than 0.007 overall on any data set: by 0.004 for the
world model on analgesics and by 0.007 for LightGBM on frozen entrees. Elsewhere
it leaves the error unchanged or raises it, by as much as 0.141 for the world model on soft drinks and
0.235 for LightGBM on analgesics.

\begin{table*}[t]
\centering\footnotesize\setlength{\tabcolsep}{3pt}
\caption{RQ4, transferability. WRMSSE overall / item by store of each tuned model trained on its own data set and
trained on it together with other data sets (each Dominick's category with the other five and Frat; Frat with the
six Dominick's categories; M5 with Favorita and Favorita with M5). The last row carries M5's tuned settings
unchanged to Favorita, to be read against Favorita's own-data row.}
\label{tab:rq-retailer}
\resizebox{\textwidth}{!}{\begin{tabular}{l cc cc}
\toprule
 & \multicolumn{2}{c}{World model} & \multicolumn{2}{c}{LightGBM}\\
\cmidrule(lr){2-3}\cmidrule(lr){4-5}
Data set & Own data & Pooled & Own data & Pooled\\
\midrule
Analgesics & 0.695 / 0.669 & 0.691 / 0.671 & 0.701 / 0.657 & 0.936 / 0.723 \\
Bottled juice & 0.573 / 0.615 & 0.613 / 0.610 & 0.576 / 0.626 & 0.582 / 0.653 \\
Cheese & 0.462 / 0.639 & 0.462 / 0.632 & 0.512 / 0.618 & 0.573 / 0.650 \\
Frozen entrees & 0.351 / 0.422 & 0.356 / 0.431 & 0.335 / 0.396 & 0.328 / 0.406 \\
Soft drinks & 0.744 / 0.600 & 0.885 / 0.747 & 0.721 / 0.570 & 0.805 / 0.594 \\
Snacks & 1.009 / 0.579 & 1.038 / 0.572 & 0.865 / 0.565 & 0.918 / 0.610 \\
Frat & 0.490 / 0.568 & 0.492 / 0.570 & 0.494 / 0.563 & 0.548 / 0.626 \\
M5 & 0.709 / 0.838 & 0.733 / 0.840 & 0.697 / 0.858 & 0.851 / 0.876 \\
Favorita & 0.428 / 0.690 & 0.448 / 0.699 & 0.448 / 0.722 & 0.459 / 0.738 \\
\midrule
Favorita, M5's settings & 0.447 / 0.692 & & 0.450 / 0.723 & \\
\bottomrule
\end{tabular}}
\end{table*}


\subsection{The broad view}
\label{sec:rq-broad}

Brought together, the four answers give one recommendation and one reason for
it. One forecast made at item by store and summed to every grain, as the
equal-weight average of a tuned world model and tuned gradient-boosted trees,
is the most accurate choice at the grains and on the kinds of line studied
here: it is lowest at store grain and above on eight of nine data sets, at the
item levels on all nine, and on the fast lines and the slow tail on most. The
two models on their own are close on items and stores they were trained on, and
each leads on different data sets. They part company on items and stores with
little or no history, where the world model holds its accuracy and LightGBM
under-forecasts. That is what the world model brings to the average, and where
it fits in the planning cycle: new lines, range changes and new stores, where a
retailer has the least history and the most to lose from a forecast that is too
low. These findings are for the price and promotion levers, the forecast windows
and the data sets studied; Section~\ref{sec:limits} sets out what they do not
cover.
